\section{Involution and ledger}
\label{sec:involution_and_ledger}

This section introduces the functional-equation involution, its odd
readout, and the two zero ledgers.  Everything here is elementary; the
point is to place the critical line in the pattern of an involution,
its fixed locus, and a readout that changes sign.

\begin{definition}[Functional-equation involution \(\JL\)]
\label{def:rh:fe-involution}
The functional-equation involution is the map
\(\JL:\mathbb C\to\mathbb C\),
\[
  \JL(s)=1-\bar s .
\]
If \(s=x+iy\), then \(\JL(s)=(1-x)+iy\): the map is reflection in the
vertical line \(\Real s=\frac12\).
\end{definition}

Directly from the formula, \(\JL(\JL(s))=1-\overline{1-\bar s}=s\), and
\(\JL(s)=s\) if and only if \(\Real s=\frac12\).  Thus \(\JL\) is an
involution whose fixed locus is the critical line.  The completed zeta
function satisfies \(\Lzeta(1-s)=\Lzeta(s)\) and
\(\Lzeta(\bar s)=\overline{\Lzeta(s)}\)
\citep{TitchmarshHeathBrown1986,Edwards1974}, so \(\JL\) maps zeros of
\(\Lzeta\) to zeros of \(\Lzeta\), preserving the strip
\(0<\Real s<1\).

\begin{definition}[Anti-invariant readout \(\psimRH\)]
\label{def:rh:psi-minus-rh}
The anti-invariant readout of \(\JL\) is the signed horizontal
displacement from the critical line,
\[
  \psimRH(s)=\Real s-\frac12 .
\]
\end{definition}

Since \(\Real\JL(s)=1-\Real s\), the readout is odd:
\(\psimRH(\JL(s))=-\psimRH(s)\).  It vanishes exactly on the fixed locus
of \(\JL\), so it is separating in the sense of the companion paper.

\begin{definition}[Nontrivial-zero ledger]
\label{def:rh:nontrivial-zero-ledger}
The \emph{nontrivial zeros} are the zeros \(\rho\) of \(\Lzeta\) with
\(0<\Real\rho<1\); they are exactly the zeros of \(\zeta\) in that
strip.  A \emph{nontrivial-zero ledger} is the set \(\Znt\) of
nontrivial zeros together with positive integer weights
\(m_\rho\) \((\rho\in\Znt)\).  The natural choice of \(m_\rho\) is the
multiplicity of \(\rho\) as a zero of \(\Lzeta\); every result below holds for arbitrary positive integer weights, and the formal development uses
unit weights where a choice is needed.
\end{definition}

\begin{definition}[Anti-invariant zero ledger]
\label{def:rh:anti-invariant-zero-ledger}
The \emph{anti-invariant zero ledger} is the extended nonnegative sum
\[
  \AZ=\sum_{\rho\in\Znt}m_\rho\,\bigl|\psimRH(\rho)\bigr|^2
     =\sum_{\rho\in\Znt}m_\rho\Bigl(\Real\rho-\frac12\Bigr)^2
  \in[0,\infty],
\]
defined as the supremum of its finite partial sums.
\end{definition}

No convergence assumption is needed: the value \(+\infty\) is allowed,
and it is only the vanishing of \(\AZ\) that matters.  Because every
summand is nonnegative and every weight is positive, \(\AZ=0\) exactly
when every nontrivial zero lies on the critical line;
\Cref{sec:translation_theorem} states this as Theorem~T.  For a finite
set \(W\subseteq\Znt\), the \emph{window sum}
\(A_W=\sum_{\rho\in W}m_\rho(\Real\rho-\frac12)^2\) is a real number
bounded by \(\AZ\); windows are the subject of
\Cref{sec:finite_windows}.
